
In this article, we propose \teraheapName{}, a system that eliminates \teraheapSD and \teraheapGC overheads for a large portion of the data in managed big data analytics frameworks.
\teraheapName{} extends the Java virtual machine (JVM) to use a second, high-capacity heap (H2) over a fast storage device that coexists alongside the regular heap (H1).
It eliminates \teraheapSD by providing direct access to objects in H2 and reduces \teraheapGC by avoiding costly \teraheapGC scans over objects in H2.
Frameworks use \teraheapName{} through its hint-based interface without modifications to the applications that run on top of them.
\teraheapName{} addresses three main challenges, as follows.

\paragraph{\textbf{Identifying candidate objects for H2}}
%
Big data frameworks move specific objects outside the managed heap on off-heap storage.
For instance, Spark moves off-heap intermediate results; Giraph moves the vertices and edges of the graph and the messages sent between vertices.
Frameworks organize such data (partitions) as groups of objects with a single-entry root reference~\cite{mahmud:survey}.
\teraheapName{} provides a hint-based interface that uses \emph{key-object opportunism}~\cite{hayes:key_opport} and enables frameworks to mark objects and indicate when to move them to H2.
During \teraheapGC, \teraheapName{} dynamically identifies the transitive closure of marked objects and move them to H2.

\paragraph{\textbf{Eliminating \teraheapGC cost for H2}}
%
\teraheapName{} presents a unified heap with the aggregate capacity of H1 and H2, where \emph{scans over H2 during \teraheapGC are eliminated,} to avoid expensive device I/O.
%
To achieve this, \teraheapName{} organizes H2 into regions with similar-lifetime objects and deals differently with liveness analysis and space reclamation.
%
We exploit the observation that analytics frameworks~\cite{zaharia:spark, giraph} organize groups of objects with similar lifetimes, such as cached partitions in Spark, in data structures with a single-entry root reference.
%
For liveness analysis, \teraheapName{} identifies live H2 regions by tracking forward (H1 to H2) and cross-region (into H2) references during \teraheapGC.
To identify live objects in H1, \teraheapName{} explicitly tracks backward references (H2 to H1) and fences \teraheapGC scans in H2.
\teraheapName{} tracks backward references using a card table optimized for storage-backed heaps, minimizing I/O traffic to the underlying device during \teraheapGC.
%
For space reclamation, the collector reclaims H1 objects as usual.
For H2 regions, unlike existing region-based allocators~\cite{gog:broom,khanh:yak} \teraheapName{} resolves the space-performance trade-off for reclaiming space differently.
Existing allocators reclaim region space eagerly by moving live objects to another region, which would generate excessive I/O for storage-backed regions.
Instead, \teraheapName{} uses the high capacity of NVMe SSDs to reclaim entire regions lazily, avoiding slow object compaction on the storage device.

\paragraph{\textbf{Applying \teraheapName{}}}
%
Managed big data analytics frameworks exhibit significant diversity in handling dataset sizes exceeding memory limits.
%
One line of frameworks, such as Spark and Giraph, reduce memory pressure by temporarily offloading excess data off-heap over a storage device.
Spark users explicitly store immutable cached data on the device, while Giraph transparently (without user hints) offloads mutable objects to the device.
We modify the two frameworks to use \teraheapName{}.
The use of \teraheapName{} is different in each framework: Spark uses \teraheapName{} to store immutable intermediate results, whereas Giraph uses \teraheapName{} to store mutable objects, such as edges and messages.
%
Another set of frameworks, such as the Neo4j Graph Data Science (GDS) library is built on the assumption that all data fits in DRAM and does not move data off-heap.
To process large datasets that exceed DRAM capacity in GDS, we show how it can exploit \teraheapName{} to offload objects to the storage device without refactoring its source code.
%Managed big data analytics frameworks exhibit significant diversity
%concerning the objects they move off-heap. We investigate how Spark
%and Giraph, two widely used frameworks, resolve the trade-off between
%GC cost due to large heaps and the overhead of off-heap accesses.
%Spark users explicitly store immutable cached data on the device,
%while Giraph transparently (without user hints) offloads mutable
%objects to the device. We modify the two frameworks to use \teraheapName{}.
%The use of \teraheapName{} is different in each framework: Spark uses \teraheapName{}
%to store immutable intermediate results, whereas Giraph uses \teraheapName{}
%to store mutable objects, such as edges and messages.

We implement \teraheapName{} and its mechanisms in OpenJDK8 and OpenJDK17, extending the Parallel Scavenge garbage collector.
We also extend the interpreter and the C1 and C2 Just-in-Time (JIT) compilers to support object updates in H2 during application execution.
%
Our evaluation shows that \teraheapName{} improves performance by up to 73\%, 28\%, and 26\% compared to native Spark, Giraph, and GDS, respectively.
%
\teraheapName{} provides similar or better performance by consuming up to $4.6\times$ and $1.2\times$ less DRAM capacity than native Spark and Giraph, respectively.
Also, it outperforms Panthera~\cite{wang:Panthera}, a garbage collector specialized for hybrid memories by up to 69\%.

Overall, our work makes the following contributions:

\begin{itemize}
	\item
	      We introduce a dual heap approach to reduce \teraheapSD and memory pressure in big data frameworks by adding a second, high-capacity, managed heap over a fast storage device.
	\item
	      We propose a hint-based interface based on key-object opportunism that enables frameworks to mark candidate objects in a coarse-grain manner and select when to move them to the second heap.
	\item
	      We show the applicability of \teraheapName{} as: (1) a large, on-heap, compute cache in Spark to store intermediate results, (2) a high-capacity heap in Giraph to store messages and edges, and (3) a high-capacity heap in GDS to offload objects during graph processing.
\end{itemize}

This article builds upon the \teraheapName{} system~\cite{teraheap} by presenting several extensions and improvements.
%
First, we rigorously analyze the trade-off between a large managed heap to store long-lived objects versus moving these objects off-heap.
%
Second, we extend the Parallel Scavenge (PS) garbage collector in OpenJDK17 to work with \teraheapName{} and provide an in-depth evaluation of its performance, showing that \teraheapName{} still provides better performance in newer JVM versions.
The key challenge in supporting OpenJDK17 includes multithreaded compactions during major \teraheapGC{} in the H1.
%
Third, we illustrate the applicability of \teraheapName{} in enabling in-memory frameworks that do not move objects off-heap to handle larger problem sizes (beyond DRAM) without needing \textit{ad-hoc} solutions.
For this purpose, we show how \teraheapName{} applies to and improves the performance of the GDS library compared to mapping the managed heap or part of the managed heap over a storage device.
%
Fourth, we investigate how dividing a fixed amount of DRAM between H1 and (OS) page cache for H2 affects \teraheapName{}'s performance.
Our goal is to understand better the scenarios in which the size of H1 or the size of the page cache has a more significant impact on performance.
%
Fifth, we examine the sensitivity of \teraheapName{} to different storage devices technologies (predominant practice in the data center) by allocating H2 over SATA SSD and NVMe devices.
%
Finally, we introduce a taxonomy of prior art and discuss the critical limitations of the state of the art in depth.

